\section{Композиция состояния: кто пишет состояние}
\label{sec:composition}

Разделы~\ref{sec:capability}--\ref{sec:economics} показали, что ущерб идёт от
отброшенной истории и что агент расплачивается за него, заново добывая те же
сведения через свои инструменты. Отсюда напрашивается лечение: если состояние
уже хранит факты, которые агент вечно выводит заново, добывать их повторно
не понадобится.

Этот раздел варьирует одну вещь: какую часть состояния \emph{рантайм} пишет
сам, не спрашивая модель. Это ось, отличная от хвоста, и мы описываем её
отдельно, потому что она отвечает на другой вопрос. Ось хвоста спрашивает,
сколько истории хранить. Эта ось спрашивает: при фиксированном объёме
истории (никакой) — кто должен писать состояние, которое её заменяет.

% Figure: the third runtime - state written by the RUNTIME, not by the model.
%
% Same visual language as fig1_architecture (box column, group frame, growth
% panel) so the three designs can be compared at a glance. The point of the
% figure is WHERE the state block comes from: in Figure 1 the model emits it as
% generated text, here the runtime derives it from what already happened.
%
% Contains no measured numbers; the measured ladder is the composition figure.

\begin{figure*}[t]
\centering
\resizebox{\textwidth}{!}{%
\begin{tikzpicture}[
  font=\small,
  box/.style={draw, rounded corners=2pt, minimum width=3.4cm,
              minimum height=0.6cm, align=center, inner sep=2pt},
  hist/.style={box, fill=red!8},
  state/.style={box, fill=blue!10, draw=blue!60, thick},
  rtstate/.style={box, fill=green!12, draw=green!55!black, thick},
  plain/.style={box, fill=black!3},
  llm/.style={box, fill=black!8, minimum height=0.85cm, thick},
  rt/.style={box, fill=green!8, draw=green!55!black, minimum height=0.85cm, thick},
  grp/.style={draw, dashed, rounded corners=3pt, inner sep=7pt},
  arr/.style={-latex, thick},
]

\node[font=\bfseries] at (0,4.55) {Состояние, написанное моделью (\SkillState{})};
\node[plain]  (p1) at (0, 3.05) {спецификация задачи $P$};
\node[state]  (s1) at (0, 2.35) {состояние выполнения $\Sigma_{t-1}$};
\node[plain]  (o1) at (0, 1.65) {последнее наблюдение $o_t$};
\node[grp, fit=(p1)(s1)(o1)] (g1) {};
\node[above=2pt of g1, font=\bfseries\footnotesize, text=blue!60!black]
  {промпт ограничен};
\node[llm] (m1) at (0,0.35) {языковая модель};
\draw[arr] (g1.south) -- (m1.north);
\node[plain, dashed] (a1) at (0,-0.85) {выполнить действие $a_t$};
\draw[arr] (m1) -- (a1);
\draw[arr, blue!60!black] (m1.west) -- ++(-1.3,0) |- (s1.west)
  node[pos=0.26, left, font=\scriptsize, text=blue!60!black, align=center]
  {модель выдаёт\\$\Delta\Sigma_t$ как текст};
\node[font=\scriptsize, text=red!65!black, align=center] at (0,-1.75)
  {состояние стоит выходных токенов;\\модель может перезаписать свои факты};

\begin{scope}[xshift=7.8cm]
\node[font=\bfseries] at (0,4.55) {Состояние, написанное рантаймом (эта работа)};
\node[plain]   (p2) at (0, 3.05) {спецификация задачи $P$};
\node[rtstate] (s2) at (0, 2.35) {проверенное состояние $V_{t-1}$};
\node[plain]   (o2) at (0, 1.65) {последнее наблюдение $o_t$};
\node[grp, fit=(p2)(s2)(o2)] (g2) {};
\node[above=2pt of g2, font=\bfseries\footnotesize, text=green!45!black]
  {промпт ограничен};
\node[llm] (m2) at (0,0.35) {языковая модель};
\draw[arr] (g2.south) -- (m2.north);
\node[plain, dashed] (a2) at (0,-0.85) {выполнить действие $a_t$};
\draw[arr] (m2) -- (a2);

\node[rt] (rec) at (4.6,-0.85) {регистратор рантайма};
\draw[arr, green!50!black] (a2.east) -- (rec.west)
  node[midway, below=1pt, font=\scriptsize, text=green!45!black, align=center]
  {результат\\инструмента};
\draw[arr, green!50!black] (rec.north) -- ++(0,1.45) -| (s2.east)
  node[pos=0.28, right, font=\scriptsize, text=green!45!black, align=left]
  {выведено, а не\\сгенерировано};
\node[font=\scriptsize, text=green!35!black, align=center] at (0,-1.75)
  {состояние не стоит выходных токенов;\\модель не может его перезаписать};
\end{scope}

\draw[dashed, black!30] (-3.4,-2.45) -- (11.2,-2.45);

\node[font=\bfseries\small, align=center] at (3.9,-3.05)
  {Что рантайм может записать, не спрашивая модель};
\begin{scope}[yshift=-3.95cm]
  \foreach \x/\tag/\what in {
     -3.0/v1/{затронутые файлы},
     -0.6/v2/{+ содержимое},
      1.8/v3/{+ команды},
      4.2/v4/{+ диагностика},
      6.6/v5/{+ окружение},
      9.0/v6/{+ тесты},
     11.4/v7/{+ провалы}} {
    \node[draw, rounded corners=2pt, fill=green!10, draw=green!55!black,
          minimum width=2.05cm, minimum height=0.62cm, font=\scriptsize,
          align=center] at (\x,0) {\textbf{\tag}\\\what};
  }
  \draw[arr, green!50!black] (-4.2,-0.68) -- (12.5,-0.68)
    node[midway, below, font=\scriptsize, text=green!45!black]
    {больше делегировано рантайму, один компонент на ступень};
\end{scope}

\end{tikzpicture}}
\caption{Два способа ограничить промпт блоком состояния, различающихся тем,
кто его пишет. \textbf{Слева:} модель выдаёт обновление состояния как
сгенерированный текст, так что его поддержание тарифицируется по выходному
тарифу, а модель вольна перезаписать ранее записанный факт, чему оригинал
приписывает \SI{68}{\percent} своих отказов моделей с открытыми
весами\origin. \textbf{Справа:} рантайм выводит ту же запись из результатов
инструментов, которые уже произошли. Блок — это факт о том, что случилось, а
не утверждение модели об этом, он не стоит сгенерированных токенов и не
может быть перезаписан моделью. Полоса внизу — это фактор, варьируемый в
\S\ref{sec:composition}: каждая ступень передаёт рантайму ещё один
компонент, при замороженных хвосте и схеме, так что разрыв между соседями
называет один делегированный компонент.}
\label{fig:runtime-written}
\end{figure*}


\subsection{Фактор}

В \SkillState{} блок состояния — сгенерированный текст: модель выдаёт
$\Delta\Sigma_t$ в той же генерации, что и действие
(Алгоритм~\ref{alg:state}). Из \S\ref{sec:economics} следуют две причины,
почему это дорого. Во-первых, поддержание состояния тарифицируется по
выходному тарифу $\beta$ — самому дорогому классу в Ур.~\eqref{eq:cost}.
Во-вторых, модель вольна перезаписать факт, который сама записала раньше, и
собственная таксономия ошибок оригинала относит на счёт именно этой ошибки
\SI{68}{\percent} отказов моделей с открытыми весами\origin.

Альтернатива — поручить рантайму фиксировать то, что уже произошло.
Результат работы инструмента — наблюдённый факт, а не утверждение: рантайм
видел, какие файлы тронуты, какая команда выполнена, что она напечатала и
прошли ли тесты. Запись этих фактов в состояние не требует генерации, и
модель не может их перезаписать.

Сколько записывать — это и есть фактор. Лестница добавляет по одному
компоненту на ступень:

\begin{center}
\begin{tabular}{@{}llp{0.52\linewidth}@{}}
\toprule
Ступень & Добавляет & Что пишет рантайм \\
\midrule
v1 & файлы        & пути, к которым обращался агент \\
v2 & + содержимое    & записи изменений по файлам и глобальные \\
v3 & + команды   & нормализованные записи команд с их выводом \\
v4 & + диагностика& структурированный вывод компилятора и линтера \\
v5 & + окружение& версии, пути, состояние интерпретатора \\
v6 & + тесты   & вызовы тестов и их результаты \\
v7 & + провалы   & записи о том, что пробовали и что не сработало \\
\bottomrule
\end{tabular}
\end{center}

Хвост заморожен на $\tailk{=}0$, схема — на \texttt{L2-semantic} в каждой
ячейке, так что ни одна ступень не заимствует преимущество у факторов из
\S\ref{sec:capability}. Соседние ступени различаются ровно одним компонентом,
и этот инвариант проверяется по флагам, которые условия реально выдают
перед запуском ячейки, а не по замыслу\code.

\paragraph{Контроль выполняется внутри этой серии.} Мы пытались повторно
использовать контроль E1, и данные это отвергли: та же рука на тех же
\VTasks{} задачах набрала 0.793 в E1 и 0.552 здесь, разойдясь по 31 из
них\live. Стенд дрейфует между сериями, поэтому заимствованный контроль
измеряет не только условие, но и этот дрейф. Против заимствованного контроля
первая ступень выглядела как дефицит 46--2; против контроля внутри серии —
как 25--2. Отсюда правило: один контроль на серию.

\subsection{Решаемость растёт по мере делегирования}

\begin{table}[t]
\centering
\caption{Лестница композиции. Стоимость, ходы и вывод — медианы по задачам,
которые решили и ступень, и контроль; отношения ниже $1$ в пользу ступени.}
\label{tab:ladder}
\input{tables/v1_ladder}
\end{table}

% The composition ladder: capability and cost against how much is delegated.
% Coordinates come from figures/data/ladder.dat, written by
% analysis/build_composition.py from the run log.

\begin{figure}[t]
\centering
\begin{tikzpicture}[font=\small]
\begin{axis}[
  name=L,
  width=0.92\linewidth, height=4.6cm,
  ybar, bar width=17pt,
  ymin=0, ymax=0.68,
  ylabel={доля решённых},
  xtick={0,1,2,3},
  xticklabels={v1 файлы,v3 +команды,v6 +тесты,v7 +провалы},
  xticklabel style={font=\scriptsize},
  enlarge x limits=0.17,
  grid=major, grid style={black!12},
  clip=false,
]
\addplot[draw=black!60, fill=green!30] table[x=x, y=rate]
  {figures/data/ladder.dat};
\draw[dashed, thick, red!70!black]
  ({rel axis cs:0,0}|-{axis cs:0,\VBaseRate}) --
  ({rel axis cs:1,0}|-{axis cs:0,\VBaseRate});
\node[font=\scriptsize, text=red!70!black, anchor=south east]
  at ({rel axis cs:0.99,0}|-{axis cs:0,\VBaseRate}) {контроль-транскрипт};
\end{axis}

\begin{axis}[
  name=C, at={($(L.south)-(0,1.65cm)$)}, anchor=north,
  width=0.92\linewidth, height=4.0cm,
  ymin=0.6, ymax=2.6,
  ylabel={стоимость, отношение к контролю},
  xtick={0,1,2,3},
  xticklabels={v1,v3,v6,v7},
  xmin=-0.4, xmax=3.4,
  grid=major, grid style={black!12},
  nodes near coords={\scriptsize$\pgfmathprintnumber[fixed,precision=2]{\pgfplotspointmeta}\times$},
  point meta=y,
  every node near coord/.append style={anchor=south west, xshift=1pt},
]
\addplot[only marks, mark=*, mark size=2.6pt, draw=black!70, fill=blue!45]
  table[x=x, y=cost] {figures/data/ladder.dat};
\draw[dashed, black!45] (axis cs:-0.4,1) -- (axis cs:3.4,1);
\node[font=\scriptsize, text=black!55, anchor=north east]
  at (axis cs:3.35,0.99) {паритет с контролем};
\end{axis}
\end{tikzpicture}
\caption{Делегирование большей части состояния рантайму двигает способность
решать задачи и стоимость в одном направлении, что противоположно обычному
компромиссу. \textbf{Сверху:} доля решённых задач по ступени против
контрольной группы транскрипта внутри серии (пунктир). Первая ступень
восстанавливает \VRateVOne{} задач; передача команд, тестов и записей о
провалах достигает \VRateVSix, закрывая примерно половину дистанции до
контрольной группы без восстановления какого-либо транскрипта.
\textbf{Снизу:} медианная стоимость на задачах, которые решили и ступень, и
контрольная группа. Те же добавления, что поднимают способность решать
задачи, также снижают стоимость с \VCostVOne$\times$ контрольной группы до
\VCostVSix$\times$, ниже паритета. Состояние, выведенное рантаймом, дешевле
состояния, написанного моделью, потому что оно никогда не генерируется по
выходному тарифу.}
\label{fig:ladder}
\end{figure}


Таблица~\ref{tab:ladder} приводит лестницу. Первая ступень — состояние,
хранящее только список затронутых файлов, — решает \VRateVOne{} против
\VBaseRate{} у контроля. Добавление записей команд поднимает долю до
\VRateVThree, добавление тестов — до \VRateVSix. На последней ступени тренд
ломается: добавление записей о провалах даёт \VRateVSeven, на одну задачу
меньше, чем v6, и при \VStepVSixVSevenLo{} против \VStepVSixVSevenHi{}
несовпадающих задач это неотличимо от отсутствия изменений. На вершине
лестница не монотонна, хотя спад остаётся внутри полосы шума и не переходит
порог парного сравнения.

Парные шаги приведены в Таблице~\ref{tab:vsteps}. Сам по себе значим только
первый шаг: переход от файлов к командам выигрывает \VStepVOneVThreeHi{}
задач и проигрывает \VStepVOneVThreeLo{} ($p={}$\VStepVOneVThreeP). Более
поздние отдельные шаги при таком размере выборки неотличимы от шума
(\VStepVThreeVSixLo{} против \VStepVThreeVSixHi{}, $p={}$\VStepVThreeVSixP;
и \VStepVSixVSevenLo{} против \VStepVSixVSevenHi,
$p={}$\VStepVSixVSevenP). А вот сквозной шаг значим: переход от v1 к v7
выигрывает \VStepVOneVSevenHi{} и проигрывает \VStepVOneVSevenLo{}
($p={}$\VStepVOneVSevenP).

Поэтому мы утверждаем тренд, а не отдельные ступени. Данные подтверждают:
состояние, которое несёт больше фактов, выведенных рантаймом, решает больше
задач, чем состояние, несущее их меньше, и основная часть выигрыша приходит
рано — там, где в запись входят команды и их вывод. Чтобы локализовать
оставшиеся компоненты относительно друг друга, нужны испытания, которые ещё
идут.

\begin{table}[t]
\centering
\caption{Парные шаги вдоль лестницы, одна и та же задача и испытание.}
\label{tab:vsteps}
\input{tables/v2_steps}
\end{table}

\subsection{Стоимость падает одновременно}

Результат, который имеет значение, — в последних трёх столбцах
Таблицы~\ref{tab:ladder}: решаемость и стоимость движутся вместе, а не друг
против друга.

На v1 рука стоит \VCostVOne$\times$ от контроля при равном исходе. Это тот
самый режим из \S\ref{sec:economics}: состояние, почти ничего не
записывающее, вынуждает агента заново добывать почти всё, и это стоит
\VTurnsVOne$\times$ от числа ходов контроля. К v6 то же сравнение даёт
\VCostVSix$\times$ — уже ниже контроля, а ходы падают до \VTurnsVSix$\times$.
Рука, которая делегирует больше всего, на этой оси одновременно самая
способная и самая дешёвая.

Механизм виден в столбце вывода. На v1 рука генерирует \VOutVOne$\times$
выходных токенов контроля; к v6 это уже \VOutVSix$\times$. За факты, которые
записал рантайм, не платят по выходному тарифу, поэтому перенос компонента
от модели к рантайму убирает сгенерированные токены и заодно устраняет то
повторное добывание, которое эти токены всё равно не могли предотвратить.

На доллар по всей серии порядок тот же: \VCostPerVOne{} решённых задач на
доллар на v1 против \VCostPerVSix{} на v6, при контроле \VBaseCostPer.

\subsection{Чего это не показывает}

Лестница не дотягивается до контроля. Каждая ступень остаётся позади него по
парной решаемости (Таблица~\ref{tab:vcontrol}), и даже лучшая всё ещё в
минусе на \VCtlVSixNet{} задач против агента с транскриптом. Делегировать
рантайму больше — строгое улучшение \emph{внутри состояния, написанного
моделью}, но на этом корпусе это пока не заменяет сохранение недавней
истории.

В сочетании с \S\ref{sec:capability}, где шести дословных обменов хватает,
чтобы закрыть разрыв полностью, практический вывод таков: два механизма
дополняют друг друга, а не соревнуются. Состояние, написанное рантаймом, —
дешёвый способ нести устойчивые факты; короткий хвост — дешёвый способ нести
недавний контекст. На этом корпусе ни один не заменяет другого.

\begin{table}[t]
\centering
\caption{Каждая ступень в паре с контролем транскрипта внутри серии.}
\label{tab:vcontrol}
\input{tables/v3_vs_control}
\end{table}

\paragraph{Статус этих чисел.} \VRuns{} запусков, \VTasks{} задач, одно
испытание. У E1 три испытания; у этой оси одно, а в этом проекте уже
случался эффект, значимый при двух испытаниях и умерший на третьем
(\S\ref{sec:pitfalls}). Направление согласуется на каждой ступени, сквозное
сравнение значимо, но отдельные шаги между ступенями следует читать как
предварительные, пока не придут остальные испытания. Ни одно утверждение
этого раздела не опирается на счётчики инструментирования плагина: по всей
серии они тождественно нулевые, и мы трактуем это как мёртвое
инструментирование, а не как свидетельство отсутствия\live.
